%! Graphing Quadratic Functions & Transformations — Unit 3, Lesson 3.1 — Homework (Answer Key)
\documentclass[10pt]{article}
\usepackage{algebra2-article}
\usepackage{algebra2-key}   % pulls in -boxes; do NOT also load -boxes
\newcommand{\TallMath}[1]{$\displaystyle #1\rule[-1.4em]{0pt}{3.2em}$}
\newcommand{\lgrid}[4]{%
  \draw[step=1, linegray!40, very thin] (#1,#3) grid (#2,#4);
  \draw[->, semithick, charcoal] (#1,0)--(#2,0) node[below right, font=\scriptsize]{$x$};
  \draw[->, semithick, charcoal] (0,#3)--(0,#4) node[left, font=\scriptsize]{$y$};}

\begin{document}

\pageheader{Unit 3, Lesson 3.1}{Homework --- Answer Key}

\vspace{0.10in}

\begin{notesbox}{Practice}
Show your reasoning. Read the feature straight from the form that advertises it.

\begin{enumerate}[leftmargin=1.9em, itemsep=11pt]
  \item \textbf{Transformations of $y=x^2$.} For each, describe the move(s) and give the new vertex.
    \begin{center}
    \renewcommand{\arraystretch}{1.5}
    \begin{tabularx}{\linewidth}{@{}l X l@{}}
    \textbf{Equation} & \textbf{Transformation of $y=x^2$} & \textbf{Vertex} \\
    \hline
    $y=x^2-5$      & \ans{down $5$} & \ans{$(0,-5)$} \\
    $y=(x+4)^2$    & \ans{left $4$} & \ans{$(-4,0)$} \\
    $y=-3x^2$      & \ans{reflect over $x$-axis, stretch by $3$ (narrower)} & \ans{$(0,0)$} \\
    $y=(x-2)^2+1$  & \ans{right $2$, up $1$} & \ans{$(2,1)$} \\
    \end{tabularx}
    \end{center}

  \item \textbf{Vertex form.} For $y=(x-3)^2-4$: vertex \ans{$(3,-4)$}, axis $x=$ \ans{$3$}, opens
        \ans{up}, min/max value \ans{min, $-4$}. Then expand to standard form: \ans{$x^2-6x+5$}.

  \item \textbf{Standard form + graph.} The graph is $f(x)=x^2+2x-3$.
    \begin{center}
    \begin{tikzpicture}[scale=0.5]
      \lgrid{-4}{2}{-5}{3}
      \draw[navy, dashed, semithick] (-1,-5)--(-1,3);
      \begin{scope}\clip (-4,-5) rectangle (2,3);
        \draw[very thick, forest, domain=-3.6:1.6, samples=75, smooth] plot (\x,{(\x)*(\x) + 2*(\x) - 3});
      \end{scope}
      \fill[charcoal] (-1,-4) circle (3pt);
      \fill[charcoal] (0,-3) circle (3pt);
      \fill[charcoal] (-3,0) circle (3pt);
      \fill[charcoal] (1,0) circle (3pt);
    \end{tikzpicture}
    \end{center}
    Axis $x=-\dfrac{b}{2a}=$ \ans{$-1$} \quad Vertex: \ans{$(-1,-4)$} \quad $y$-intercept: \ans{$(0,-3)$}
    \quad Range: \ans{$y\ge -4$}.

  \item \textbf{Intercept form.} For $g(x)=(x-2)(x-6)$: zeros $x=$ \ans{$2$} and $x=$ \ans{$6$};
        axis $x=$ \ans{$4$}; vertex \ans{$(4,-4)$}; $y$-intercept \ans{$(0,12)$}.

  \item \textbf{Matching.} Write the letter of the graph that matches each equation.
    \begin{center}
    \begin{tikzpicture}[scale=0.36]
      % P: (x-1)^2 - 4  vertex (1,-4) up
      \begin{scope}
        \lgrid{-2}{4}{-5}{3}
        \begin{scope}\clip (-2,-5) rectangle (4,3);
          \draw[very thick, forest, domain=-1.4:3.4, samples=60, smooth] plot (\x,{(\x-1)*(\x-1) - 4});
        \end{scope}
        \fill[charcoal] (1,-4) circle (2.5pt);
        \node[charcoal, font=\scriptsize] at (1,-6.6){\textbf{P}};
      \end{scope}
      % Q: -(x+2)^2 + 1  vertex (-2,1) down
      \begin{scope}[xshift=8cm]
        \lgrid{-5}{1}{-4}{2}
        \begin{scope}\clip (-5,-4) rectangle (1,2);
          \draw[very thick, forest, domain=-4.2:0.2, samples=60, smooth] plot (\x,{-1*(\x+2)*(\x+2) + 1});
        \end{scope}
        \fill[charcoal] (-2,1) circle (2.5pt);
        \node[charcoal, font=\scriptsize] at (-2,-5.6){\textbf{Q}};
      \end{scope}
      % R: x^2 + 2  vertex (0,2) up
      \begin{scope}[xshift=16cm]
        \lgrid{-3}{3}{0}{6}
        \begin{scope}\clip (-3,0) rectangle (3,6);
          \draw[very thick, forest, domain=-2:2, samples=60, smooth] plot (\x,{(\x)*(\x) + 2});
        \end{scope}
        \fill[charcoal] (0,2) circle (2.5pt);
        \node[charcoal, font=\scriptsize] at (0,-1.4){\textbf{R}};
      \end{scope}
    \end{tikzpicture}
    \end{center}
    $y=(x-1)^2-4:$ \ans{P} \qquad $y=-(x+2)^2+1:$ \ans{Q} \qquad $y=x^2+2:$ \ans{R}

  \item \textbf{Explain.} Which form would you choose to read each feature the fastest, and why?
        Zeros: \ans{intercept form} \; Vertex: \ans{vertex form} \; $y$-intercept: \ans{standard form}.
        \ansline{Each form displays one feature directly: intercept form sets each factor to zero for the zeros, vertex form shows $(h,k)$, and standard form gives $(0,c)$ by reading $c$.}
\end{enumerate}
\end{notesbox}

\begin{extensionbox}
\textbf{A fountain arc.} A garden fountain's water arcs so its height (in feet) at horizontal distance
$x$ feet from the nozzle is $h(x) = -x(x-6)$ (given in intercept form).
\begin{enumerate}[label=(\alph*), leftmargin=1.8em, itemsep=5pt]
  \item Where does the water leave the ground and where does it land? (the zeros)
  \item What is the maximum height, and how far out does it occur? (the vertex)
  \item At $x=1$ ft the water is $5$ ft high. Use \emph{symmetry} --- no algebra --- to name another
        distance where it is again $5$ ft high.
\end{enumerate}
\ansline{(a) Zeros $x=0$ and $x=6$: it leaves the ground at the nozzle ($x=0$) and lands $6$ ft away. \;\; (b) Axis $x=3$ (halfway between $0$ and $6$); $h(3)=-3(-3)=9$, so the maximum height is $9$ ft, $3$ ft out. \;\; (c) $x=1$ is $2$ ft left of the axis $x=3$, so its mirror is $2$ ft right: $x=5$ ft (check $h(5)=-5(-1)=5$).}
\end{extensionbox}

\begin{spiralbox}
\textbf{Coming next --- Lesson 3.2: Solving Quadratics by Factoring.} Today you \emph{read} zeros from
intercept form. Next you will learn to \emph{produce} that factored form --- turning $x^2+2x-3$ into
$(x+3)(x-1)$ --- so you can find the zeros of a quadratic given in standard form by hand.
\end{spiralbox}

\end{document}
